\documentclass[11pt, a4paper]{article}
\usepackage[a4paper, top=2.5cm, bottom=2.5cm, left=2cm, right=2cm]{geometry}
\usepackage{amsmath}
\usepackage{amssymb}
\usepackage{booktabs}
\usepackage{enumitem}

\setlength{\parindent}{0pt}
\setlength{\parskip}{1em}

\begin{document}

\begin{center}
    {\Large \textbf{Astrodynamics 2026-1: Lecture Notes}} \\
    \vspace{0.2cm}
    {\huge \textbf{Finite Maneuvers \& Low-Thrust Spirals}} \\
    \vspace{0.4cm}
    \textit{Prof. Juan | Universidad de Antioquia - Aerospace Engineering}
\end{center}

\hrule
\vspace{0.5cm}

\section*{1. Introduction to Continuous Thrust}
Throughout the earlier modules, we modeled orbit transfers (like the Hohmann transfer) using the \textbf{Impulsive Burn} assumption. We assumed that engines produce massive amounts of thrust over negligible time frames, meaning the position vector ($\boldsymbol{r}$) remains constant while the velocity vector ($\boldsymbol{v}$) changes instantaneously.

However, as detailed in \textit{Vallado (Section 6.7)}, modern interplanetary missions frequently utilize Solar Electric Propulsion (SEP) systems, such as ion or Hall-effect thrusters. These engines provide microscopic thrust (often on the order of millinewtons) but can fire continuously for months or years. We must now model \textbf{Finite Maneuvers}, where the spacecraft's position, velocity, and mass change continuously and simultaneously.

\section*{2. The Equations of Motion (GMAT Math Specifications)}
When analyzing continuous thrust, we can no longer rely on purely analytical geometric solutions (like Patched Conics). We must define the exact Ordinary Differential Equations (ODEs) and rely on numerical integrators (like the Runge-Kutta solvers in GMAT) to propagate the state vector.

As outlined in the GMAT Mathematical Specifications, the total acceleration acting on a spacecraft is the sum of gravitational forces, environmental perturbations, and the active engine thrust:

\begin{equation}
    \ddot{\boldsymbol{r}} = -\frac{\mu}{r^3}\boldsymbol{r} + \boldsymbol{a}_{perturb} + \frac{\boldsymbol{T}}{m(t)}
\end{equation}

Where $\boldsymbol{T}$ is the thrust vector produced by the engine. Crucially, the mass $m(t)$ is not constant. As the engine fires, it depletes propellant according to the mass flow rate equation:

\begin{equation}
    \dot{m} = -\frac{|\boldsymbol{T}|}{I_{sp} g_0}
\end{equation}

Where $I_{sp}$ is the Specific Impulse of the engine and $g_0$ is standard gravity ($9.80665 \text{ m/s}^2$). Because mass is continuously decreasing, a constant thrust vector $\boldsymbol{T}$ will result in a steadily \textit{increasing} acceleration profile over the duration of the mission.

\section*{3. Analytical Approximation: The Low-Thrust Spiral}
While numerical integration is required for precise targeting, mission designers need analytical approximations for initial mission sizing. 

If a low-thrust engine fires continuously in the tangential direction (parallel to the velocity vector), and the thrust-to-weight ratio is extremely small ($T/W < 10^{-4}$), the orbit's eccentricity remains near zero. The orbit behaves as a \textit{quasi-circular spiral} slowly unwinding away from the central body.

Vallado demonstrates that the total change in velocity ($\Delta V$) required to spiral from an initial circular orbit of radius $r_1$ to a final circular orbit of radius $r_2$ can be approximated as the absolute difference between the circular velocities at those two altitudes:

\begin{equation}
    \Delta V_{spiral} \approx |v_{c1} - v_{c2}| = \left| \sqrt{\frac{\mu}{r_1}} - \sqrt{\frac{\mu}{r_2}} \right|
\end{equation}

\subsection*{3.1 The Gravity Penalty}
If we calculate the $\Delta V$ for a Hohmann transfer between the same two orbits ($r_1$ to $r_2$), we will always find that:
\begin{equation}
    \Delta V_{spiral} > \Delta V_{Hohmann}
\end{equation}

This discrepancy is known as \textbf{gravity loss}. Impulsive maneuvers are efficient because they apply $\Delta V$ exactly at periapsis (maximizing the Oberth effect). Continuous thrust systems waste energy by thrusting at inefficient points in the orbit.

\subsection*{3.2 Why Use Low-Thrust?}
If the spiral transfer requires more $\Delta V$, why do we use it? The answer lies in the Rocket Equation:
\begin{equation}
    \Delta m_{propellant} = m_0 \left( 1 - e^{-\frac{\Delta V}{I_{sp} g_0}} \right)
\end{equation}

Chemical rockets typically possess an $I_{sp}$ of $\sim 300$ seconds. Electric thrusters possess an $I_{sp}$ of $\sim 3,000$ seconds. Even though the required $\Delta V$ is higher, the exponential efficiency of the electric thruster results in massive propellant savings, allowing for much larger scientific payloads. \textbf{We trade transfer time for payload mass.}

\section*{4. Transitioning to GMAT Operations}
In your upcoming GMAT labs, you will transition from using the \texttt{Target} and \texttt{ImpulsiveBurn} commands. Instead, you will build physical representations of the spacecraft hardware (\texttt{ChemicalThruster}, \texttt{ElectricThruster}, and \texttt{FuelTank}). 

You will use the \texttt{BeginFiniteBurn} and \texttt{EndFiniteBurn} commands wrapped around a \texttt{Propagate} sequence. GMAT will use the ODEs defined above to rigorously integrate the spiral trajectory, naturally accounting for gravity losses, mass depletion, and even thrust termination during eclipses (when solar panels cannot power the thrusters).

\end{document}